\begin{figure}[!htbp]
\centering
\begin{tikzpicture}[>=Latex,x=1cm,y=1cm,every node/.style={font=\small}]
 \draw[->] (0,0.45)--(0,-4.45);
 \node[anchor=east] at (-.2,.35) {root};
 \foreach \y/\label in {-.65/{parent cut},-1.85/{child cut},-3.05/{old sample level},-4.1/{old descendants}} {
  \draw (-.08,\y)--(.08,\y);
  \node[anchor=east] at (-.2,\y) {\label};
 }
 \node[anchor=west] at (.25,-.65) {$m x_p\le(1-\eta)m$};
 \node[anchor=west] at (.25,-1.85) {$m x_c\ge(1-\eta)m$};
 \node[anchor=west] at (.25,-3.05) {$m$};
 \node[anchor=west] at (.25,-4.1) {$t_2,\ldots,t_d$};
 \draw[decorate,decoration={brace,amplitude=5pt}] (4.3,-1.65)--(4.3,-3.25);
 \node[anchor=west,align=left,text width=3.1cm] at (4.6,-2.45)
  {new processed block:\\exponents lie in $[m-\eta,m]$};
\end{tikzpicture}
\caption{One structural transition when both cuts are proper. The parent
exponent decreases by a fixed factor; the levels strictly below it through
the old sample level form the new processed block. The cut parameters are
$x_p=\Prob(Q<q)$ and $x_c=\Prob(Q<r)$. Trivial parent or child cuts are
omitted, as specified in the proof. The old descendant row is absent when
$d=1$.}
\label{fig:hidden-cuts}
\end{figure}
